\Needspace{14\baselineskip}
\subsection{Density roots from a compact K\"ahler factor}
\label{subsec:closed-factor-root-spaces}

On the products of \cref{subsec:closed-factor-product-densities} the lower mass vanishes, so $\mathcal D_H$ is a proper subspace of $L^N$ by \cref{prop:individual-density-banach}, and we do not know whether it equals $\mathcal D_H^0$. Nevertheless, \cref{thm:closed-factor-product-equation} computes the root functional exactly on all roots that come from the K\"ahler factor, and there it is given by the K\"ahler formula of the same proposition. We keep the notation of that subsection: $(Y,\kappa)$ is a connected compact K\"ahler manifold of dimension $k$ with $\int_Y\kappa^k=1$, the product $X=X_0\times Y$ has dimension $N$ and background $H=h_0+\kappa$, with $\int_XH^N=CV_0$, and $R=\mu/\mathrm{d}V$ has minimum $a$ and maximum $A$. The root functional $\rho_H$ and the spaces $\mathcal D_H^0\subset\mathcal D_H\subset L^N(\mathrm{d}V)$ are those of \cref{sec:integrable-gauduchon-tests,subsec:bounded-root-approximation}, on the $N$-dimensional manifold $X$. For a real function $\xi\in L^N(Y,\kappa^k)$ put
\[
 J\xi=R^{1/N}\pi_Y^*\xi.
\]

\Needspace{14\baselineskip}
\begin{proposition}\label{prop:closed-factor-root-subspace}\index{Density space!Kahler factor subspace@K\"ahler factor subspace}
For every real $\xi\in L^N(Y,\kappa^k)$,
\[
 \|J\xi\|_{L^N(\mathrm{d}V)}=\|\xi\|_{L^N(\kappa^k)},\qquad
 \rho_H(J\xi)=(CV_0)^{-1/N}\|\xi\|_{L^N(\kappa^k)}.
\]
Thus $J$ maps $L^N(Y,\kappa^k)$ linearly onto a closed subspace of $\mathcal D_H^0$, with $\|J\xi\|_{\mathcal D_H}=(1+(CV_0)^{-1/N})\|\xi\|_{L^N(\kappa^k)}$, and every $p$ in this subspace satisfies
\[
 \rho_H(p)=\Bigl(\int_XH^N\Bigr)^{-1/N}\|p\|_{L^N(\mathrm{d}V)}.
\]
The tails satisfy, for every $K>0$,
\[
 \rho_H\bigl(|J\xi|\mathbf1_{\{|J\xi|>K\}}\bigr)
 \leq\Bigl(\frac1{CV_0}\int_{\{|\xi|^N>K^N/A\}}|\xi|^N\kappa^k\Bigr)^{1/N},
\]
with equality, and with $A=1$, if $\mathrm{d}V=\mu$. Finally, every real $p\in L^N(\mathrm{d}V)$ with $|p|\leq\pi_Y^*G$ for some $G\in L^N(Y,\kappa^k)$ belongs to $\mathcal D_H^0$, and $\rho_H(p)\leq(aCV_0)^{-1/N}\|G\|_{L^N(\kappa^k)}$.
\end{proposition}
\begin{proof}
Put $g=|\xi|^N$, so that $|J\xi|^N=Rg$. By Fubini's theorem, $\int_X|J\xi|^N\,\mathrm{d}V=\int_Xg\,\mu=\int_Yg\kappa^k$, which is the first identity. For $\xi\ne0$, \cref{thm:closed-factor-product-equation} gives $b_H(|J\xi|^N)=CV_0/\|\xi\|_{L^N(\kappa^k)}^N$, which is the second identity; for $\xi=0$ both sides vanish. Since $J$ is linear, the two identities applied to $\xi-\eta$ give
\[
 \|J\xi-J\eta\|_{\mathcal D_H}=(1+(CV_0)^{-1/N})\|\xi-\eta\|_{L^N(\kappa^k)}.
\]
Hence the image of $J$ is complete, and therefore closed in $\mathcal D_H$. The formula for $\rho_H$ on the image combines the two identities with $\int_XH^N=CV_0$. Bounded functions are dense in $L^N(Y,\kappa^k)$ and have bounded images, so the image of $J$ lies in the closure of the bounded functions, which is $\mathcal D_H^0$.

On $\{|J\xi|>K\}$ we have $A|\xi|^N\geq R|\xi|^N>K^N$. Hence
\[
 |J\xi|\mathbf1_{\{|J\xi|>K\}}\leq J\bigl(|\xi|\mathbf1_{\{|\xi|^N>K^N/A\}}\bigr),
\]
and the monotonicity of $\rho_H$ (\cref{prop:individual-root-triangle}) together with the second identity gives the tail bound. If $\mathrm{d}V=\mu$, then $R=1$ and $J\xi=\pi_Y^*\xi$, and the two sides of the last display coincide. Finally, if $|p|\leq\pi_Y^*G$, then $|p|\leq a^{-1/N}J|G|$. Since $\mathcal D_H$ is an order ideal in $L^N$ and $\mathcal D_H^0$ is an order ideal in $\mathcal D_H$ (\cref{prop:individual-density-banach,prop:root-uniform-tails}), $p$ belongs to $\mathcal D_H^0$, and monotonicity gives the bound for $\rho_H(p)$.
\end{proof}

On this subspace the root functional is thus the same as on a K\"ahler manifold of total volume $\int_XH^N$, and the uniform tails are controlled by the $L^N$ tails of $\xi$ alone. Since $R$ is bounded above and below and $\xi$ is arbitrary, $\mathcal D_H^0$ contains functions in no $L^r$ with $r>N$. For the density of the explicit equation in \cref{thm:zero-lower-mass-product-equation} the rate is logarithmic.

\Needspace{12\baselineskip}
\begin{example}\label{ex:closed-factor-logarithmic-tails}\index{Uniform tails!logarithmic decay}
Let $Y=\mathbb P^k$ with the normalized Fubini--Study form, let $\mathrm{d}V=\mu$ and $0<\alpha<1$, and let $q_\alpha$ be the density of \cref{lem:projective-sublogarithmic-density}. The root $p_\alpha=\pi_Y^*q_\alpha^{1/N}$ of the density in \cref{thm:zero-lower-mass-product-equation} belongs to $\mathcal D_H^0$, belongs to no $L^r(\mathrm{d}V)$ with $r>N$, and
\[
 \lim_{K\to\infty}(\log K)^{k(1-\alpha)/N}\rho_H\bigl(p_\alpha\mathbf1_{\{p_\alpha>K\}}\bigr)
 =\Bigl(\frac{\alpha^k}{CV_0}\Bigl(\frac kN\Bigr)^{k(1-\alpha)}\Bigr)^{1/N}.
\]
\end{example}
\begin{proof}
By \cref{prop:closed-factor-root-subspace} with $\xi=q_\alpha^{1/N}$, the root $p_\alpha$ belongs to $\mathcal D_H^0$, and
\[
 \rho_H\bigl(p_\alpha\mathbf1_{\{p_\alpha>K\}}\bigr)^N=\frac{T(K^N)}{CV_0},
 \qquad
 T(t)=\int_{\{q_\alpha>t\}}q_\alpha\kappa^k.
\]
By Fubini's theorem, $\int_Xp_\alpha^r\,\mathrm{d}V=\int_Yq_\alpha^{r/N}\kappa^k$, which is infinite for $r>N$ by \cref{lem:projective-sublogarithmic-density}. It remains to show that
\begin{equation}\label{eq:closed-factor-tail-mass}
 T(t)\sim\alpha^k\Bigl(\frac{\log t}{k}\Bigr)^{-k(1-\alpha)}\qquad(t\to\infty);
\end{equation}
the substitution $t=K^N$ then gives the limit.

From now on we use the notation of the proof of \cref{lem:projective-sublogarithmic-density}: $r=\sum_{i=1}^k|z_i|^2$ on the chart $z_0=1$, $s(r)=1-\log(r/(1+r))\geq1$, and the function $Q(r)$. Put $\beta=1+k(1-\alpha)$. The radial integration in that proof shows that the measure $q_\alpha\kappa^k$ gives the set $\{r<\varepsilon\}$ the mass $m(\varepsilon)=Q(\varepsilon)^k$, since $Q(0)=0$ and the point $o$ carries no mass. As $r\downarrow0$, we have $s(r)=-\log r+O(1)\to+\infty$, and the formulas there give $Q(r)\sim\alpha s(r)^{\alpha-1}$ and $Q'(r)\sim\alpha(1-\alpha)r^{-1}s(r)^{\alpha-2}$. Hence
\[
 m(\varepsilon)\sim\alpha^k(-\log\varepsilon)^{-k(1-\alpha)}\quad(\varepsilon\downarrow0),
 \qquad
 q_\alpha(r)\sim\alpha^k(1-\alpha)\varphi(r)\quad(r\downarrow0),
\]
where $\varphi(r)=r^{-k}s(r)^{-\beta}$. Since $s'(r)=-1/(r(1+r))$, the logarithmic derivative of $\varphi$ is $r^{-1}\bigl(\beta/((1+r)s(r))-k\bigr)$, which is negative for small $r$. Choose $r_0\in(0,1)$ such that $\varphi$ is decreasing on $(0,r_0)$ and
\[
 \tfrac12\alpha^k(1-\alpha)\varphi\leq q_\alpha\leq2\alpha^k(1-\alpha)\varphi
 \quad\text{on }\{0<r<r_0\},
\]
and let $M_0$ be the maximum of $q_\alpha$ on the compact set $Y\setminus\{r<r_0\}$, which does not contain $o$.

Let $t>M_0$. Since $s\geq1$, every point of $\{q_\alpha>t\}\setminus\{o\}$ satisfies $r<r_0$ and $2\alpha^k(1-\alpha)r^{-k}>t$, so $\{q_\alpha>t\}\setminus\{o\}$ lies in $\{r<\varepsilon_+(t)\}$, where $\varepsilon_+(t)=(2\alpha^k(1-\alpha)/t)^{1/k}$, and $T(t)\leq m(\varepsilon_+(t))$. Conversely, put $\varepsilon_-(t)=t^{-1/k}(\log t)^{-2\beta/k}$. For large $t$, $\varepsilon_-(t)<r_0$ and $s(\varepsilon_-(t))\leq2\log t$, so for $0<r<\varepsilon_-(t)$,
\[
 q_\alpha(r)\geq\tfrac12\alpha^k(1-\alpha)\varphi(\varepsilon_-(t))
 \geq\tfrac12\alpha^k(1-\alpha)\,t(\log t)^{2\beta}(2\log t)^{-\beta}>t
\]
for large $t$. Hence $T(t)\geq m(\varepsilon_-(t))$. Since $-\log\varepsilon_\pm(t)=k^{-1}\log t+O(\log\log t)$, the asymptotics of $m$ give \eqref{eq:closed-factor-tail-mass} from both sides.
\end{proof}

These results concern the roots that come from the K\"ahler factor and the roots they dominate. They do not decide whether $\mathcal D_H^0=\mathcal D_H$: a root with nonvanishing uniform tails would have to escape every bound $|p|\leq\pi_Y^*G$ with $G\in L^N(Y,\kappa^k)$. Nor do they show that uniform root tails suffice for an exact equation. The exact equations of \cref{thm:closed-factor-product-equation} come from the product structure. For general densities that depend on both factors and belong to no $L^p$ with $p>1$, the existence of an exact equation remains open even when their roots are dominated by such a bound.
